% ===============================================
% 별표5 국문초록  (2쪽 이내)
% ===============================================
\clearpage
\thispagestyle{plain}

\begin{center}
{\fontsize{16pt}{22pt}\selectfont\bfseries 국문초록\par}

\vspace{8mm}

{\fontsize{14pt}{20pt}\selectfont\bfseries \ThesisTitleKR\par}
\end{center}

\vspace{6mm}

\begin{flushright}
{\fontsize{12pt}{18pt}\selectfont
\AuthorKRPlain\par
\DepartmentKR\par
\GraduateSchoolKR\par}
\end{flushright}

\vspace{8mm}

% 본문 — 11pt, 줄간격 180~200%, 문단 시작마다 2칸 들여쓰기
\setstretch{1.9}

여기에 국문초록 본문을 입력하세요. 본문은 11\,pt, 휴먼명조 계열, 줄간격
180\,\%\,$\sim$\,200\,\%, 문단 시작마다 2칸 들여쓰기를 적용합니다. 본 템플릿은
가천대학교 일반대학원의 ``국문 학위논문 작성지침''을 근거로 작성되었으며,
이공·자연계열 번호 체계(제1장 → 1.1 → 1.1.1 → (1) → ①)를 사용합니다.

국문초록은 연구의 배경과 필요성, 목적, 방법, 주요 결과, 결론을 2쪽
이내로 간결하게 기술합니다. 본문 마지막에는 핵심어를 7$\sim$8단어 이내로
기입합니다.

\vfill

\noindent\textbf{핵심어}: 키워드1, 키워드2, 키워드3, 키워드4, 키워드5, 키워드6, 키워드7

\clearpage
